\documentclass[11pt]{article}

\usepackage[a4paper,margin=1in]{geometry}
\usepackage{amsmath,amssymb,amsthm,mathtools}
\usepackage{hyperref}
\usepackage{enumitem}
\usepackage{xcolor}
\usepackage{mdframed}

\newtheorem{problem}{Problem}
\theoremstyle{definition}
\newtheorem{definition}{Definition}
\theoremstyle{remark}
\newtheorem*{why}{Why it matters}

\newcommand{\bbN}{\mathbb{N}}
\newcommand{\bbZ}{\mathbb{Z}}
\newcommand{\src}[1]{\par\smallskip\noindent{\footnotesize\textsf{Source: #1}}}
\newcommand{\bp}[1]{\begin{mdframed}[backgroundcolor=black!4,linewidth=0.4pt]#1\end{mdframed}}

\title{An Understanding Kit\\[2pt]\large Eight results from the corpus, to rebuild by hand}
\author{Patrick S.\ Mahon}
\date{October 2026}

\begin{document}
\maketitle

\bp{\textbf{How to use this.} Read one problem's statement. Close the PDF and
work on paper until you see \emph{why} it must be true. That means rebuilding
it, not verifying it. Hints are in Part~II and proofs in Part~III, deliberately
far away. The problems run from a few minutes to an hour. None needs anything
outside this document. A problem that lights up tells you something; one that
doesn't tells you something too.}

\tableofcontents

\part{Problems}

\section*{Shared definitions}
\addcontentsline{toc}{section}{Shared definitions}

\begin{definition}[two-valued state]
For a family $\mathcal D$ of subsets of a set $\Omega$ closed under complements
and countable disjoint unions, a \emph{two-valued state} is $s:\mathcal
D\to\{0,1\}$ with $s(\Omega)=1$ that is additive on disjoint pairs. It is
\emph{$\sigma$-additive} if additivity holds for countable disjoint families, and
\emph{finitely additive} if only for finite ones. The \emph{Dirac} $\delta_\omega$
is $s(E)=1\iff\omega\in E$.
\end{definition}

\begin{definition}[language, two tokens]
A \emph{language} is a relation $\rho\subseteq A\times A$ on a finite alphabet $A$,
read as a digraph. It is \emph{primitive} if some power of its adjacency matrix is
entrywise positive (equivalently: strongly connected, and the gcd of its cycle
lengths is $1$). \emph{Two tokens} sit on distinct vertices and move
\emph{synchronously}: at each step each token crosses one arc of $\rho$. They may
never occupy the same vertex at the same time.
\end{definition}

% ------------------------------------------------------------------
\section{Tiny}

\begin{problem}[stationary two-valued states are fixed points]\label{p:stat}
Let $A$ be countable and $\Omega=A^{\bbZ}$ with the product $\sigma$-algebra.
\begin{enumerate}[label=(\alph*)]
\item Every $\sigma$-additive two-valued state on $\Omega$ is a Dirac
  $\delta_\omega$.
\item If it is also shift-invariant ($s(S^{-1}E)=s(E)$), then $\omega$ is a
  constant sequence.
\end{enumerate}
\end{problem}
\begin{why}
``Two-valued'' means \emph{one trajectory}, not ``deterministic dynamics''. So a
genuine stationary process is never two-valued. This is why the $\sigma$-essential
lane sits off the path to dynamics.
\end{why}
\src{not written elsewhere; derived 2026-10-09 in conversation.}

\begin{problem}[out-of-phase tokens never collide]\label{p:phase}
Let $\lambda:R\to\bbZ_c$ be a function on a set of vertices. Token~1 is only ever
at $w\in R$ at times $t\equiv\lambda(w)$, and token~2 only at times
$t\equiv\lambda(w)+e \pmod c$, with $e\not\equiv0$. Show they never collide in~$R$.
\end{problem}
\begin{why}
This is the engine of the reach-lemma attack. A phase difference is a
collision-proof shield.
\end{why}
\src{\texttt{papers/reconstruction/notes/joint2\_wielandt\_finding.md}, ``(2)
Out-of-phase tokens never collide''.}

\begin{problem}[the ascending ring]\label{p:ring}
Fix $L\ge2$ and let $C_k=\{x\in\bbN^L : x_k<x_{k+1}\}$, indices mod $L$.
\begin{enumerate}[label=(\alph*)]
\item No point $x$ lies in every $C_k$.
\item Suppose that, for each window $\{k,k+1\}$, $s_k$ is a $\sigma$-additive
  two-valued state on the $\sigma$-algebra of events depending only on
  $(x_k,x_{k+1})$. Suppose also that adjacent windows agree on events of their
  shared coordinate, and that $s_k(C_k)=1$ for all $k$. Show this is impossible.
\end{enumerate}
\end{problem}
\begin{why}
Drop ``$\sigma$-'' in (b) and the configuration \emph{exists}: an ordered product
ultrafilter makes every window an ascent. That finitely additive object, glued on a
carrier, is the Lean-checked countable $\sigma$-essential witness (every $L\ge4$).
\end{why}
\src{\texttt{notes/open\_questions/delay\_embedding/collision\_vs\_generation.md}
(update 2026-10-09); Lean \texttt{formalization/QuerySystem/QuerySystem/WindowGluing.lean},
\texttt{sL\_sigmaEssential}.}

% ------------------------------------------------------------------
\section{Small}

\begin{problem}[the cycle lemma]\label{p:cycle}
Let $\rho$ be primitive and $C=(x_0,\dots,x_{c-1})$ a simple cycle, $c\ge2$. Two
tokens \emph{chase} on $C$ at offset $j\neq0$: both follow $C$, token~2 being $j$
positions ahead. Show that:
\begin{enumerate}[label=(\alph*)]
\item from any offset, by legal non-colliding moves and back again, the tokens can
  reach a chase at any other offset;
\item among the closed two-token runs so produced there are lengths with gcd $1$.
\end{enumerate}
\end{problem}
\begin{why}
This is the reason chases are \emph{aperiodic}. Aperiodicity is exactly what the
flagship's forcing lemma needs.
\end{why}
\src{\texttt{joint2\_wielandt\_finding.md}, ``CYCLE LEMMA (proved
2026-10-09)''.}

\begin{problem}[one chase class]\label{p:kstar}
Let $C,C'$ be simple cycles of length $\ge2$ sharing a vertex $v$. Find a single
synchronous step that turns some chase on $C$ into a chase on $C'$. Conclude, using
that the simple cycles of a strongly connected digraph are chained by shared
vertices, that all chases on all cycles are mutually reachable.
\end{problem}
\src{\texttt{joint2\_wielandt\_finding.md}, ``LEMMA K*''.}

\begin{problem}[no witness in the Boolean case]\label{p:bool}
Let $\mathcal L$ be a Boolean $\sigma$-algebra on $\Omega$ and $B\subseteq\mathcal L$
finite and closed under complement. Let $s_0$ be a two-valued assignment on $B$
that extends to \emph{some} finitely additive two-valued state $\mu$ on $\mathcal L$.
Show $s_0$ extends to a Dirac.
\end{problem}
\begin{why}
This is the baseline the whole $\sigma$-essential lane pushes against. Coherence
supplies $\mu$, and distributivity turns $\mu$ into a point. A witness must break
the second step.
\end{why}
\src{\texttt{papers/sigma\_essential/sigma\_essential\_body.tex}, Prop.~\texttt{prop:boolean}.}

% ------------------------------------------------------------------
\section{Medium}

\begin{problem}[Dirac test and localization]\label{p:local}
With $B$ finite and complement-closed and $s_0$ two-valued on $B$, put
$K(s_0)=\bigcap\{E\in B: s_0(E)=1\}$.
\begin{enumerate}[label=(\alph*)]
\item $\delta_\omega$ extends $s_0$ iff $\omega\in K(s_0)$.
\item Hence a pattern that is finitely coherent but has no $\sigma$-additive
  extension needs $K(s_0)=\varnothing$ \emph{and} must kill every non-Dirac
  $\sigma$-additive extension. Explain why the second clause is where all the
  difficulty lives.
\end{enumerate}
\end{problem}
\src{\texttt{sigma\_essential\_body.tex}, Lemma \texttt{lem:dirac-iff}, Thm
\texttt{thm:localization}, Remark \texttt{rem:localization-scope}.}

\begin{problem}[the parity theorem, odd half]\label{p:parity}
The golden-mean language on $\{0,1\}$ forbids $11$. Observe it around a ring of
length $L$ through the $L$ adjacent-pair windows. Let each window assign
probability $\tfrac12$ to each of $01$ and $10$.
\begin{enumerate}[label=(\alph*)]
\item The windows agree on shared sites.
\item For odd $L$, no probability measure on legal ring configurations has these
  window marginals.
\item \emph{Stretch.} For even $L$ every coherent datum is realisable. (Hint
  section only.)
\end{enumerate}
\end{problem}
\begin{why}
This is the flagship's smallest real instance of contextuality, and it is
determined by parity.
\end{why}
\src{\texttt{papers/reconstruction/notes/archive/reconstruction\_body\_PROVED\_ARCHIVE.tex},
\S``The parity theorem''.}

\bigskip
\noindent\textbf{For comparison.} Re-read your own CE / Stone route
(\texttt{papers/paper\_i/probability\_from\_observation\_body.tex}, ``Stone
realization from finitely additive data'') and notice what the click felt like.

% ==================================================================
\newpage
\part{Hints}

\begin{enumerate}[label=\textbf{\arabic*.}]
\item (a) For each time $k$, the sets $\{x_k=a\}$, $a\in A$, are a countable
  partition. How many can be true? (b) What is $s\circ S^{-1}$ for a Dirac?
\item Compare the two congruences at a hypothetical meeting.
\item (a) Follow the inequalities around. (b) Show each $s_k$ is a Dirac at some
  pair, then glue.
\item An \emph{excursion} is a walk $x_i\to x_{i'}$ whose interior avoids $C$.
  While one token takes it, the other stays on $C$. Track the offset. Which
  quantity changes it? Why do these quantities generate $\bbZ_c$? (Think about
  closed walks at $x_0$ and primitivity.) Removing one vertex from a Cayley graph:
  what keeps it connected?
\item Put token~1 at $v$ and token~2 one step behind $v$ on $C$.
\item The finite Boolean algebra generated by $B$ has atoms. Exactly one of them
  is true.
\item (a) Use complements: every false constraint is a true constraint in disguise.
\item (b) If every window is a.s.\ $01$ or $10$, what must a global configuration
  look like? (c) The coherence polytope is the fractional stable-set polytope of
  the cycle. Look at $\det(I+P)$ for the cyclic shift $P$.
\end{enumerate}

% ==================================================================
\newpage
\part{Proofs}

\paragraph{Problem~\ref{p:stat}.}
(a) For each $k$ the events $\{x_k=a\}$ ($a\in A$) are countably many, disjoint,
and cover $\Omega$, so by $\sigma$-additivity exactly one, say $a_k$, has value
$1$. Let $\omega=(a_k)_k$. Then $\{\omega\}=\bigcap_k\{x_k=a_k\}$, a countable
intersection of value-$1$ events. Its complement is a countable union of
value-$0$ events, which can be disjointified into value-$0$ events, so it has
value $0$ and $s(\{\omega\})=1$. Hence $s(E)=1\iff\omega\in E$.
(b) $\delta_\omega\circ S^{-1}=\delta_{S\omega}$. Invariance gives
$\delta_{S\omega}=\delta_\omega$, so $S\omega=\omega$, so $\omega$ is constant.
\hfill$\square$

\paragraph{Problem~\ref{p:phase}.}
A meeting at $w$ at time $t$ needs $t\equiv\lambda(w)$ and
$t\equiv\lambda(w)+e$, so $e\equiv0$, a contradiction.\hfill$\square$

\paragraph{Problem~\ref{p:ring}.}
(a) $x_0<x_1<\dots<x_{L-1}<x_0$. Equivalently, take $k$ maximising $x_k$; then
$x_k<x_{k+1}$ contradicts maximality.
(b) As in Problem~\ref{p:stat}(a), each $s_k$ is $\delta_{(a_k,b_k)}$ on its
window, since $\bbN^2$ is countable and the singletons are measurable. Agreement
on the shared coordinate gives $b_k=a_{k+1}$. So one point $x=(a_k)$ makes every
$C_k$ true, which contradicts (a).\hfill$\square$

\paragraph{Problem~\ref{p:cycle}.}
\emph{Case $c=2$.} The single offset is the swap $(x_0,x_1)\to(x_1,x_0)$, a
closed run of length~$1$.

\emph{Case $c\ge3$.}
\begin{enumerate}[label=(\roman*)]
\item \emph{Defect.} An excursion $x_i\to x_{i'}$ of length $\ell$ has
  \emph{defect} $\Delta=\ell-((i'-i)\bmod c)$. Loops and chords count.
\item \emph{Moves.} If token~1 takes it while token~2 stays on $C$, then token~2
  advances $\ell$ while token~1 lands $\ell-\Delta$ further along $C$. So the
  offset changes by $+\Delta$, and by $-\Delta$ if token~2 moves. The interior
  avoids $C$, so the only possible collision is at the landing. The move is
  therefore legal iff the new offset is $\neq0$.
\item \emph{The defects generate $\bbZ_c$.} A closed walk at $x_0$ splits into
  arcs of $C$ and excursions, so its length is $\equiv\sum\Delta\pmod c$. By
  primitivity these lengths have gcd $1$, so they generate $\bbZ_c$.
\item \emph{Connectivity, giving (a).} The offset graph is
  $\mathrm{Cay}(\bbZ_c,\pm\{\Delta\})$ with the vertex $0$ removed. A connected
  vertex-transitive graph of degree $\ge2$ is $2$-connected (Watkins 1970), so it
  stays connected.
\item \emph{Lengths, giving (b).} The chase itself is a closed run of length $c$.
  For each $\Delta$, start at an offset $\notin\{0,-\Delta\}$. Token~1 takes the
  excursion, then token~2 takes it when it arrives. That is a closed run of length
  $\equiv\Delta$. Since the $\Delta$'s generate $\bbZ_c$, the gcd is $1$.
  \hfill$\square$
\end{enumerate}

\paragraph{Problem~\ref{p:kstar}.}
Token~1 is at $v$ and token~2 at $v$'s predecessor on $C$. In one step token~1
moves to $v$'s successor on $C'$ (which is $\ne v$, since $|C'|\ge2$) and token~2
moves onto $v$. Both are now on $C'$, one apart: a chase on $C'$. The symmetric
step goes back. Chains of cycles sharing vertices connect any two cycles, and a
loop in a chain can be skipped, because its neighbours share its vertex. With
Problem~\ref{p:cycle}(a), everything is one class.\hfill$\square$

\paragraph{Problem~\ref{p:bool}.}
The finite algebra $B_0$ generated by $B$ has atoms that partition $\Omega$. By
finite additivity exactly one atom $a$ has $\mu(a)=1$, and $a\ne\varnothing$.
Every element of $B_0$ contains $a$ or is disjoint from it. So for $\omega\in a$,
$\delta_\omega$ agrees with $\mu$, hence with $s_0$, on $B$.\hfill$\square$

\paragraph{Problem~\ref{p:local}.}
(a) $\delta_\omega$ extends $s_0$ iff $\omega\in E$ for every true $E$, and
$\omega\notin E$ for every false $E$. A false $E$ has true complement in $B$, so
its constraint is already a true one. Hence the condition is just
$\omega\in K(s_0)$.
(b) $\sigma$-additive two-valued states are Dirac or not. (a) handles the Diracs.
The empty-kernel condition is easy to arrange (Navara--Pt\'ak does it with even a
$\sigma$-additive global extension). So the whole difficulty is destroying every
non-Dirac $\sigma$-additive repair while keeping finite coherence.
\hfill$\square$

\paragraph{Problem~\ref{p:parity}.}
(a) Every window gives each site value $1$ with probability $\tfrac12$.
(b) Almost surely every window is $01$ or $10$, so a realising configuration
alternates $0,1,0,1,\dots$ around the ring. That needs $L$ even.
(c) The coherence polytope is $\{u\in[0,1]^L: u_i+u_{i+1}\le1\}$ and $R$ is the
stable-set polytope. The constraint matrix is $I+P$, with
$\det(I+P)=1-(-1)^L$; its proper square submatrices are path-incidence matrices,
with determinants in $\{0,\pm1\}$. For even $L$ the matrix is totally unimodular,
so the polytope is integral and $C=R$. Full proof: archive
\S``The parity theorem''.\hfill$\square$

% ==================================================================
\newpage
\appendix
\part{Appendix: definitions, notation and background results}

Everything the problems use, in one place. The table says what to read for each
problem. Entries marked \emph{(standard)} are textbook facts, quoted without
proof.

\medskip
\begin{center}
\begin{tabular}{cl}
\textbf{Problem} & \textbf{Read} \\ \hline
1 & A.1--A.4, A.6 \\
2 & C.6 \\
3 & A.1--A.4, A.7 \\
4 & B.1--B.6, C.1--C.5, D.1--D.3 \\
5 & B.1--B.4, B.7, C.1--C.3 \\
6 & A.1--A.5, E.1--E.3 \\
7 & A.1--A.5, E.1--E.5 \\
8 & F.1--F.5 \\
\end{tabular}
\end{center}

% ------------------------------------------------------------------
\section{Measure-theoretic basics}

\paragraph{A.1 $\sigma$-algebra, Dynkin system.} A \emph{$\sigma$-algebra} on
$\Omega$ is a family of subsets containing $\Omega$ and closed under complements
and countable unions. A \emph{Dynkin system} (or $\sigma$-class) only requires
closure under complements and \emph{countable disjoint} unions. It need not be
closed under intersections. That gap is exactly what makes non-Boolean
(``quantum-logic-like'') carriers possible.

\paragraph{A.2 Product $\sigma$-algebra.} On $\Omega=A^{I}$ (sequences indexed by
$I$, e.g.\ $I=\bbZ$ or $\{0,\dots,L-1\}$) it is the smallest $\sigma$-algebra
making each coordinate $x\mapsto x_k$ measurable. It is generated by the
\emph{cylinders} $\{x: x_{k_1}=a_1,\dots,x_{k_m}=a_m\}$. For countable $A$, every
singleton $\{\omega\}$ of $A^{\bbZ}$ is measurable, being a countable intersection
of cylinders.

\paragraph{A.3 Two-valued state; finite vs.\ $\sigma$-additivity.} See the shared
definitions in Part~I. For a two-valued $s$:
\begin{itemize}[nosep]
\item finite additivity means: of two disjoint sets, at most one is true, and
  their union is true iff one of them is;
\item $\sigma$-additivity means the same for countable disjoint families.
\end{itemize}
Consequences:
\begin{itemize}[nosep]
\item $s(E^c)=1-s(E)$;
\item $s$ is monotone ($E\subseteq F$, $s(E)=1\Rightarrow s(F)=1$);
\item in a $\sigma$-algebra, a countable intersection of true sets is true when
  $s$ is $\sigma$-additive.
\end{itemize}

\paragraph{A.4 Dirac.} $\delta_\omega(E)=1\iff\omega\in E$. It is always
$\sigma$-additive. ``Two-valued $\sigma$-additive $=$ Dirac'' holds on any
$\sigma$-algebra in which a countable partition into singletons (or a countable
separating family) is available. Problem~\ref{p:stat}(a) is an instance.

\paragraph{A.5 Boolean algebra, atoms.} A \emph{Boolean algebra} of sets is closed
under complement, finite union and finite intersection. A \emph{finite} one is
generated by finitely many sets, and its \emph{atoms} (minimal nonempty members)
partition $\Omega$. Every member is a union of atoms
\emph{(standard)}. ``Boolean'' $=$ distributive $=$ intersection-closed, in the
concrete setting.

\paragraph{A.6 Shift, stationarity.} The shift $S:A^{\bbZ}\to A^{\bbZ}$ is
$(Sx)_k=x_{k+1}$. A state $s$ is \emph{stationary} (\emph{shift-invariant}) if
$s(S^{-1}E)=s(E)$ for all $E$. For Diracs, $\delta_\omega\circ
S^{-1}=\delta_{S\omega}$. A \emph{process} is a stationary probability measure on
$A^{\bbZ}$.

\paragraph{A.7 Ultrafilter (for the ``why it matters'' notes).} An
\emph{ultrafilter} $p$ on $\bbN$ is a finitely additive two-valued state on all
subsets of $\bbN$. It is \emph{non-principal} if it is not a Dirac, i.e.\ every
finite set is false. Non-principal ultrafilters exist by Zorn's lemma, and they
are \emph{never} $\sigma$-additive. The \emph{ordered product} $p\otimes p$ makes
$\{(a,b): a<b\}$ true: ``the second coordinate is chosen after, and above, the
first''. This is the device that fakes an ascending ring.

% ------------------------------------------------------------------
\section{Digraphs and languages}

\paragraph{B.1 Digraph vocabulary.}
\begin{itemize}[nosep]
\item A \emph{walk} is a sequence of vertices joined by arcs. Its
  \emph{length} is the number of arcs.
\item A walk is \emph{closed} if it ends where it starts.
\item A \emph{simple cycle} is a closed walk with no repeated vertex other than
  the endpoints. A \emph{loop} is an arc $v\to v$, a cycle of length $1$.
\end{itemize}

\paragraph{B.2 Strongly connected.} Every vertex reaches every other by a walk.

\paragraph{B.3 Period, primitive.} For a strongly connected digraph, the
\emph{period} is the gcd of all closed-walk lengths, equivalently of all
simple-cycle lengths. It is \emph{aperiodic} (\emph{primitive}) if the period is
$1$. Equivalently, some power of the adjacency matrix is entrywise positive, and
then all larger powers are too \emph{(standard; Perron--Frobenius / Wielandt)}.
In words, for large $n$ there is a walk of length exactly $n$ between any two
vertices.

\paragraph{B.4 Language.} A relation $\rho\subseteq A\times A$, used as the
digraph of legal consecutive symbols. ``Primitive language'' means its digraph is
primitive.

\paragraph{B.5 Excursion, chord, defect.} Fix a simple cycle
$C=(x_0,\dots,x_{c-1})$.
\begin{itemize}[nosep]
\item An \emph{excursion} is a walk from some $x_i$ to some $x_{i'}$ whose
  interior vertices are all off $C$, other than one of $C$'s own arcs.
\item A \emph{chord} is an excursion of length $1$, i.e.\ an arc $x_i\to x_{i'}$
  that is not $C$'s arc.
\item Its \emph{defect} is $\Delta=\ell-((i'-i)\bmod c)$: how much longer the
  excursion is than the arc of $C$ it bypasses, mod $c$.
\end{itemize}

\paragraph{B.6 Decomposing a closed walk.} A closed walk starting and ending on $C$
splits uniquely into arcs of $C$ and excursions, by cutting at every visit to
$C$.

\paragraph{B.7 Chain-connectivity of cycles \emph{(standard; proved in the repo,
session 5)}.} In a strongly connected digraph, any two simple cycles are joined by
a chain of simple cycles in which consecutive cycles share a vertex.

% ------------------------------------------------------------------
\section{Two tokens}

\paragraph{C.1 Pair graph $D$.} Vertices are ordered pairs $(a,b)$ with $a\neq b$
(the two tokens' positions). There is an arc $(a,b)\to(c,d)$ iff $a\to c$ and
$b\to d$ are arcs of $\rho$ and $c\neq d$. A walk in $D$ is a synchronous
two-token run that never collides.

\paragraph{C.2 Unordered pair graph $\bar D$.} The same, with $(a,b)$ and $(b,a)$
identified. A closed walk in $\bar D$ returns the tokens to the same \emph{set} of
positions, possibly swapped. A \emph{swap} (crossed walk) is a run
$(a,b)\leadsto(b,a)$.

\paragraph{C.3 Chase, offset.} On a simple cycle $C$, a \emph{chase} has both
tokens on $C$, following it, with token~2 a fixed number $j\neq0$ of positions
ahead (the \emph{offset}). It is a closed run of length $c$. A pair of vertices is
\emph{cocyclic} if both lie on one simple cycle, and then it is a chase position.

\paragraph{C.4 SCC, recurrent.} A \emph{strongly connected component} (SCC) is a
maximal set of mutually reachable vertices. A pair is \emph{recurrent} if it lies
on a closed walk of $D$ (equivalently, of $\bar D$), i.e.\ it can return to
itself.

\paragraph{C.5 Aperiodic SCC; why it matters.} An SCC is aperiodic if the gcd of
its closed-walk lengths is $1$. The flagship's forcing lemma needs every SCC of
$\bar D$ that carries a swap to be aperiodic. Otherwise its unsafe ring lengths
could sit in a residue class like $3\bbZ$, which would be a ``fork''.

\paragraph{C.6 Phase.} If walks from a start $y$ to $w$ all have the same length
mod $c$, call it $\lambda(w)$. A token starting at $y$ at time $0$ is at $w$ only
at times $\equiv\lambda(w)$. A second token can be \emph{out of phase} by $e$:
always at $w$ at times $\equiv\lambda(w)+e$.

% ------------------------------------------------------------------
\section{Cyclic groups and Cayley graphs}

\paragraph{D.1 $\bbZ_c$, generation.} $\bbZ_c=\{0,\dots,c-1\}$ under addition mod
$c$. A set $S$ \emph{generates} $\bbZ_c$ iff $\gcd(S\cup\{c\})=1$. If a set of
integers has gcd $1$, its residues mod $c$ generate $\bbZ_c$.

\paragraph{D.2 Cayley graph.} $\mathrm{Cay}(\bbZ_c,\pm S)$ has vertices $\bbZ_c$
and an edge $g$---$g\pm s$ for each $s\in S$. It is connected iff $S$ generates
$\bbZ_c$, and it is \emph{vertex-transitive}: every vertex looks the same.

\paragraph{D.3 Watkins' theorem \emph{(standard; Watkins 1970)}.} A connected
vertex-transitive graph of degree $d$ has vertex connectivity at least
$\tfrac23(d+1)$. In particular, for $d\ge2$ removing one vertex leaves it
connected.

% ------------------------------------------------------------------
\section{Quantum-logic side ($\sigma$-essential lane)}

\paragraph{E.1 Concrete $\sigma$-orthoposet.} A family $\mathcal L$ of subsets of
$\Omega$, containing $\Omega$, closed under complement ($E^\perp=E^c$) and under
countable unions of pairwise disjoint members. This is a Dynkin system (A.1)
viewed as a ``logic''. Disjoint sets play the role of \emph{orthogonal} events.

\paragraph{E.2 $\perp$-closed sub-orthoposet; local pattern.} A finite
$B\subseteq\mathcal L$ closed under complement: the finitely many events actually
being reported. A \emph{two-valued state on $B$} (local pattern) $s_0$ assigns
$0/1$ with $s_0(E^c)=1-s_0(E)$ (and additivity on disjoint pairs within $B$).

\paragraph{E.3 Finitely coherent.} $s_0$ extends to some \emph{finitely
additive} two-valued state on all of $\mathcal L$.

\paragraph{E.4 Kernel.} $K(s_0)=\bigcap\{E\in B: s_0(E)=1\}$, the points
compatible with every reported ``yes''.

\paragraph{E.5 $\sigma$-essential contextual state.} A local pattern that is
finitely coherent but extends to \emph{no} $\sigma$-additive two-valued state on
$\mathcal L$: coherence that only finite additivity can honour. The
Navara--Pt\'ak example (1983) has empty kernel yet a $\sigma$-additive extension,
which shows that ``empty kernel'' alone is not enough.

% ------------------------------------------------------------------
\section{Commensurability side (flagship)}

\paragraph{F.1 Ring protocol, windows.} Sites $\bbZ_L$ around a ring. Each
adjacent pair $\{i,i+1\}$ is a \emph{window}, reported as a probability table over
the legal pairs of $\rho$.

\paragraph{F.2 Coherent (EA) vs.\ realisable (PR).} The window data are
\emph{coherent} (empirically adequate) if overlapping windows give the same
single-site marginals. They are \emph{realisable} if some probability measure on
legal ring configurations has exactly these window marginals. \emph{Contextual}
means coherent but not realisable. A protocol is \emph{commensurable} if every
coherent datum is realisable.

\paragraph{F.3 Polytopes $C$, $R$.} Coherent data form a polytope $C$, and
realisable data the convex hull $R\subseteq C$ of the deterministic
configurations. Commensurable $\iff C=R$.

\paragraph{F.4 Stable sets.} In a graph, a \emph{stable (independent) set} has no
two adjacent vertices.
\begin{itemize}[nosep]
\item $\mathrm{STAB}$ is the convex hull of their indicator vectors.
\item $\mathrm{FSTAB}=\{u\in[0,1]^V: u_i+u_j\le1 \text{ on edges}\}$ is its
  fractional relaxation.
\item For the golden mean on a ring, $C\cong\mathrm{FSTAB}(C_L)$ and
  $R\cong\mathrm{STAB}(C_L)$, with $u_i=\Pr[x_i=1]$.
\end{itemize}

\paragraph{F.5 Total unimodularity \emph{(standard; Hoffman--Kruskal)}.} A matrix is
\emph{totally unimodular} if every square submatrix has determinant in
$\{0,\pm1\}$. Then $\{u: Mu\le b,\,u\ge0\}$ has integral vertices for every
integral $b$. Also used:
\begin{itemize}[nosep]
\item $\det(I+P)=\prod_k(1+\omega^k)=1-(-1)^L$ for the cyclic shift $P$, whose
  eigenvalues are the $L$-th roots of unity $\omega^k$;
\item path-incidence matrices have determinants in $\{0,\pm1\}$ (triangularise
  along the path).
\end{itemize}

\end{document}
